\newpage\handout[true]
{\textsc{Math 2106-B: Foundations of Mathematical Proof}}{Fall 2025}
{Homework 2}
{\textit{Prof.} \href{https://ceheitsch.github.io/webpage/}{Dr. Christine Heitsch}}{Due Date: September 4th}
{Math 2106-B: HW2}
{\large
  \noindent
  \textbf{Name:} \HandoutAuthor \smallskip \\

  \noindent
  \textbf{GTID:} 903743506 \smallskip \\

  \noindent
  %%% List anyone with whom you discussed the problems
  \textbf{Collaborators:} None. This material reflects independent preparation. \smallskip \\

  \noindent
  %%% List all resources used OTHER than the textbook, lecture notes, 
  %%% webwork problems, ICA questions, and previous homework assignments.
  \textbf{Outside resources:} This work was informed by MATH 2106 course materials 
  (ICAs: in-class activities and WebWork contents) from \HandoutInstitution\ (the course textbook is 
  Chartrand \parencite{chartrand_mathematical_2018}).   
  Outside references specifically include
  textbooks written by 
  Grimaldi \parencite{grimaldi_discrete_2004} (used in MATH 3012) and 
  Rosen \parencite{rosen_discrete_2019} (used in CS 2050). \smallskip \\

  \noindent
  \textbf{Notice}: The answers contain cross-references throughout this document as clickable hyperlinks in most PDF viewers. However, the Gradescope PDF view might not support clickable hyperlinks.
}
\printbibliography[
  heading=bibintoc,   
  title={References}     % change “References” text if needed
]
% `heading=bibnumbered` if you prefer it numbered chapter (in \tableofcontents)
% `heading=bibintoc`    if you prefer it unnumbered but still in the ToC\newpage
The following table presents the standard logical equivalences from these sources together with 
the author's own (or the student, myself, Jaehoon Song) clarifications and labels where applicable.


\begin{table}[h!]
  \centering
  \renewcommand{\arraystretch}{1.2}
  \setlength{\tabcolsep}{8pt}
  \begin{tabular}{|p{5.0cm}|p{5.0cm}|p{5.2cm}|}
    \hline
    \rowcolor[HTML]{E0E0E0}
    \multicolumn{2}{|l|}{\textbf{Logical Equivalences}} & \textbf{Name} \\
    \hline
    $\,p \AND q \equiv q \AND p$ & $\,p \OR q \equiv q \OR p$ & Commutative Law \\
    \hline
    $\, (p \AND q) \AND r \equiv p \AND (q \AND r)$ & $\, (p \OR q) \OR r \equiv p \OR (q \OR r)$ & Associative Law \\
    \hline
    $\,p \AND (q \OR r) \equiv (p \AND q) \OR (p \AND r)$ & $\,p \OR (q \AND r) \equiv (p \OR q) \AND (p \OR r)$ & Distributive Law \\
    \hline
    $\,\overline{p \AND q} \equiv \sNOT{p} \OR \sNOT{q}$ & $\,\overline{p \OR q} \equiv \sNOT{p} \AND \sNOT{q}$ & DeMorgan's Law \\
    \hline
    \multicolumn{2}{|l|}{$\,\overline{\overline{p}} \equiv p$} & Double Negation law \\
    \hline
    $\,p \AND \boldsymbol{T} \equiv p$ & $\,p \OR \boldsymbol{F} \equiv p$ & Identity Law \\
    \hline
    $\,p \OR \boldsymbol{T} \equiv \boldsymbol{T}$ & $\,p \AND \boldsymbol{F} \equiv \boldsymbol{F}$ & Domination Law \\
    \hline
    $\,p \AND p \equiv p$ & $\,p \OR p \equiv p$ & Idempotent Law \\
    \hline
    \multicolumn{2}{|l|}{$\,p \THEN q \equiv \sNOT{p} \OR q$} & Implication Equivalence \\
    \hline
    \multicolumn{2}{|l|}{$\,\overline{p \THEN q} \equiv p \AND \sNOT{q}$} & Negation of Implication \\
    \hline
    \multicolumn{2}{|l|}{$\,p \THEN q \equiv \sNOT{q} \THEN \sNOT{p}$} & Contrapositive Equivalence \\
    \hline
    \multicolumn{2}{|l|}{$\,p \OR \sNOT{p} \equiv \boldsymbol{T}$} & Law of Tautology \\
    \hline
    \multicolumn{2}{|l|}{$\,p \AND \sNOT{p} \equiv \boldsymbol{F}$} & Law of Contradiction \\
    \hline
  \end{tabular}
\end{table}\newpage

The following definitions are presented from standard mathematical sources (\textbf{Sets}).
Clarifications and labels have been added where applicable.

\begin{define}{Basic Set Concepts}\phantomsection\label{def:basic-set-concepts} % \hyperref[def:basic-set-concepts]{YOUR_CONTEXT}
  \begin{enumerate}
    \item \textbf{Universal Set:} A universal set $\setUNIV$ is a
    set that contains all objects under consideration in a particular context or problem.
    \item \textbf{Empty Set:} The empty set, denoted $\setNULL$,
    is the unique set that contains no elements.
    \item \textbf{Subset:} Let $A$ and $B$ be sets. $A$ is a
    subset of $B$ ($A \subseteq B$) if (and only if)
    for all $x \in A$, $x \in B$.
    \item \textbf{Set Equality:} Let $A$ and $B$ be sets. The sets $A$ and $B$ are
    equal ($A = B$) if $A \subseteq B$ and $B \subseteq A$.
    \item \textbf{Power Set:} Let $A$ be a set. The power set of $A$, denoted $\sPOWER{A}$ or $2^A$,
    is the set of all subsets of $A$. That is, $\mathcal{P}(A) = \{S\in\setUNIV : S \subseteq A\}$.
  \end{enumerate}
\end{define}
The following definitions extend the understanding of set operations
and properties.


\begin{define}{Set Operations}\phantomsection\label{def:set-operations} % \hyperref[def:set-operations]{YOUR_CONTEXT}
  \begin{enumerate}
    \item \textbf{Complementary Set:} Let $A$ be a subset of the universal set $\setUNIV$.
    The complement of $A$, denoted $A^c$ or $\overline{A}$, is the set of all elements
    in $\setUNIV$ that are not in $A$. That is, $A^c = \{x \in \setUNIV : x \notin A\}$.
    \item \textbf{Union:} Let $A$ and $B$ be sets. The union of $A$ and $B$,
    denoted $A \cup B$, is the set of all elements that are in $A$, or in $B$, or in both.
    That is, $A \cup B = \{x : x \in A \lor x \in B\}$.
    \item \textbf{Intersection:} Let $A$ and $B$ be sets. The intersection of $A$ and $B$,
    denoted $A \cap B$, is the set of all elements that are in both $A$ and $B$.
    That is, $A \cap B = \{x : x \in A \land x \in B\}$.
    \item \textbf{Set Difference:} Let $A$, $B$ be subsets of a universal set $\setUNIV$.
    The set difference $A \sMINUS B$ is defined to be the set
    $\{ x \in \setUNIV \LDIV x \in A \text{ and } x \notin B\}$. \vspace*{10px}\\
    \textit{It is sometimes called the relative complement of $B$ in $A$,
    and denoted in the textbook as $A - B$.}
    \item \textbf{Cartesian Product:} Let $A$ and $B$ be sets. The Cartesian product of
    $A$ and $B$, denoted $A \sTIMES B$, is the set consisting of all ordered pairs:
    \[
    A \sTIMES B = \{(a, b) \LDIV a \in A, b \in B\}.
    \]
  \end{enumerate}
\end{define}


\newpage

The following claims are given in the reference as well. Let $A,B\in U$, a universal set.


\begin{claim}{Subset of Union: $A \subseteq \sOR{A}{B}$}\phantomsection\label{claim:subset-union} % \hyperref[claim:subset-union]{YOUR_CONTEXT}
  \begin{subprove}[bluecolor]
      Let $x\in U$, assume $x\in A$. Then the statement \inlinequote{$x\in A$ or $x\in B$} is true.
      Hence, by the definition, $\sOR{A}{B}=\{x\in U\LDIV x\in A$ or $x\in B\}$, $x\in \sOR{A}{B}$.
      Thus, $A \subseteq \sOR{A}{B}$.
  \end{subprove}
\end{claim}

\begin{claim}{Empty Subset: $\setNULL \subseteq A$}\phantomsection\label{claim:empty-subset} % \hyperref[claim:empty-subset]{YOUR_CONTEXT}
  \begin{subprove}[bluecolor]
      Let $x\in U$, assume $x\in \setNULL$. But there are no elements in $\setNULL$, so
      $x\in \setNULL$ is false. Since the premise is false, the implication if $x\in \setNULL$,
      then $x\in A$ is true. Thus, $\setNULL \subseteq A$ by definition of \hyperref[def:basic-set-concepts]{subsets}.
  \end{subprove}
\end{claim}

\begin{claim}{Subset as Equality: $\sOR{A}{B}=A$ if and only if $B\subseteq A$}\phantomsection\label{claim:subset-equality} % \hyperref[claim:subset-equality]{YOUR_CONTEXT}
  \begin{subprove}[bluecolor]
      Let $x\in U$, first, we assume $B\subseteq A$, we have already shown
      $A \subseteq \sOR{A}{B}$ by the claim, \hyperref[claim:subset-union]{Subset of Union}. 
      To show $\sOR{A}{B} \subseteq A$, assume $x\in \sOR{A}{B}$.
      By definition, $x\in A$ or $x\in B$. Considering $x\in B$ and assumption
      $B \subseteq A$, $x\in B$ implies $x\in A$. Thus, $\sOR{A}{B} \subseteq A$.
      Hence, if $B\subseteq A$, then $\sOR{A}{B} = A$. \vspace*{10px}\\
      Now, we assume $\sOR{A}{B}=A$. To show $B\subseteq A$, we now assume $x\in B$. Since $x\in B$, by definition
      of union, $x\in \sOR{A}{B} = A$. Therefore, $B\subseteq A$.
  \end{subprove}
\end{claim}


%%%%%%%%%%%%%%%%%%% 9/8/25

Let $\setUNIV$ be a set. 
\begin{claim}{ICA6}\phantomsection\label{claim:ica6} % \hyperref[claim:ica6]{YOUR_CONTEXT}
  For all $A,B\in\setUNIV$, $$A\sTIMES B =\setNULL\THEN \group{A=\setNULL\OR B=\setNULL}$$
  \begin{subprove}[bluecolor]
    Let $A,B\in\setUNIV$. We prove $A\sTIMES B=\setNULL$. Assume it is not the case that 
    $A=\setNULL$ or $B=\setNULL$. Thus, we have $A=\setNULL$ and $B=\setNULL$. 
    Since $A\ne\setNULL$ by assumption, there exists an element $a\in A$. Likewise, there 
    is some $b\in B$ since $B\ne\setNULL$. Thus, $\group{a,b}\in\sBUILD{(x,y)}{x\in A,y\in B}$ 
    by definition. But then, $A\sTIMES B\ne\setNULL$ by definition. Hence, the claim holds. 
  \end{subprove}
\end{claim}


\begin{define}{Relatively Prime Integers (Coprime)}\phantomsection\label{def:relatively-prime} % \hyperref[def:relatively-prime]{YOUR_CONTEXT}
  Let $x,y \in \mathbb{Z}$. The integers $x$ and $y$ are said to be
  \textbf{relatively prime} if and only if
  $\gcd(x,y) = 1$.
\end{define}
\newpage
The following definitions and axioms are presented from standard mathematical sources (\textbf{Number Theory}) together with the author's own (or the student, myself, Jaehoon Song) clarifications and labels where applicable.

\begin{define}{Integers}[def:integers][red]
  $\setZ$: the set of all integers
\end{define}

\begin{define}{Rational Numbers}[def:rational-numbers][red]
  $\setQ$: the set of all rational numbers
\end{define}

\begin{define}{Real Numbers}[def:real-numbers][red]
  $\setR$: the set of all real numbers
\end{define}

\begin{define}{Natural Numbers}[def:natural-numbers][red]
  $\setN$ or $\setZ_{>0}$: the set of all natural numbers which are the set of positive integers (not including $0$)
\end{define}

\begin{define}{Positive Real Numbers}[def:positive-real-numbers][red]
  $\setR_{>0}$: the set of all positive real numbers
\end{define}

\begin{define}{Non-negative Real Numbers}[def:non-negative-real-numbers][red]
  $\setR_{\geq 0}$: the set of all non-negative real numbers
\end{define}


\newpage

The following definitions and claims extend the understanding of integers.

Here are the ICA2 Definitions and Axioms as follows.

\begin{define}{Even Integer}[def:even-integer][red]
  An integer $n$ is \underbar{even} if there exists 
  (existential quantification logic $\exists$) an integer $k$ such that
  \[
  n = 2k.
  \]
\end{define}

\begin{define}{Odd Integer}[def:odd-integer][red]
  An integer $n$ is \underbar{odd} if 
  there exists some integer $k$ such that
  \[
  n = 2k + 1.
  \]
\end{define}

\begin{define}{Closure under Addition}[def:closure-addition][red]
  The sum of two integers is again an 
  integer. (Closure under addition)
\end{define}

\begin{define}{Closure under Multiplication}[def:closure-multiplication][red]
  The product of two integers is again 
  an integer. (Closure under multiplication)
\end{define}

\begin{define}{Odd if and only if Not Even}[def:odd-iff-not-even][red]
  An integer is odd if and only if it 
  is not even.
\end{define}

Here are the ICA3 Definitions and Axioms as follows.

\begin{define}{Universal Closure under Addition}[def:universal-closure-addition][red]
  The sum of two integers (Universal quantification for all, for every, for each) is again an integer. (Closure under addition)
\end{define}

\begin{define}{Universal Closure under Multiplication}[def:universal-closure-multiplication][red]
  The product of two integers (Universal quantification for all, for every, for each) is again an integer. (Closure under multiplication)
\end{define}

\begin{define}{Even Integer Square}[def:even-integer-square][red]
  An integer $n$ is even if and only if $n^2$ is even.
\end{define}

\newpage

The following claims are given in the reference as well.

\begin{define}{Sum of Even Integers}[claim:sum-even-integers][red]
  The sum of two even integers is again an even integer.
  \begin{subprove}[red]
    Let $x$ be an integer and $y$ be an integer ($\forall x,y \in \setZ$). 
    Suppose $x$ is even and $y$ is even. Then, by \nameref{def:even-integer}, there exists an integer 
    $k$ and an integer $j$ such that $x=2k$ and $y=2j$. Then, 
    $x+y=\left(2k\right)+\left(2j\right)=2\left(k+j\right)$. By \nameref{def:closure-addition}, since $k$ and 
    $j$ are integers, then so is $\left(k+j\right)$. Hence, by definition, $x+y$ is an 
    even integer.
  \end{subprove}
\end{define}

The following definitions extend the understanding of integers including divisibility.

\begin{define}{Divisibility}[def:divisibility][red]
  An integer $n$ divides an integer $m$, denoted $n \LDIV m$, if there 
  exists an integer $k$ such that $m=n k$. Otherwise, $n \nmid m$.
\end{define}

The following definitions extend the classification of real numbers into rational and irrational numbers.

\begin{define}{Rational Number}[def:rational-number][red]
  A real number $x$ is \underbar{rational} if there exist integers $a$ and $b$ with $b \neq 0$ such that $bx = a$.
  Otherwise, $x$ is \underbar{irrational}.

  \textbf{Note}: Compare the rational number definition with the even/odd axiom \inlinequote{odd iff not even} from \nameref{def:odd-iff-not-even}.
\end{define}

The following claims are given in the reference as well.

\begin{define}{Irrationality of $\sqrt{2}$}[claim:sqrt2-irrational][red]
  $\sqrt{2}$ is irrational.
  \begin{subprove}[red]
    Suppose $\sqrt{2}$ is rational. Then, by \nameref{def:rational-number}, there exists $a,b\in \setZ$ 
    with $b\ne 0$ such that $b\cdot\sqrt{2}=a$. We may further assume that any factors 
    common to $a$ and $b$ have already been removed.\\

    Then, squaring both sides, this yields $2\cdot b^2 = a^2$.
    Since products of integers are again integers by \nameref{def:closure-multiplication}, we have that $a^2,b^2\in \setZ$.
    Thus, by definition, $a^2$ is even. Therefore, by \nameref{def:even-integer-square}, $a$ is even. 
    Thus, there exists $c\in\setZ$ such that $a=2c$. Then
    $$
    2b^2=a^2={2c}^2=4c^2
    $$
    Then, $b^2=2c^2$. Then, as before, $b^2$ is even, so $b$ is even.
    However, $a$ and $b$ both being even contradicts the fact that any factors 
    common to $a$ and $b$ have already been removed.
  \end{subprove}
\end{define}






%%%%%%%%%%%%%%%%%%% New style (refactor model)
\begin{define}{Relatively Prime Integers (Coprime)}[def:relatively-prime][red]
  Let $x,y \in \mathbb{Z}$. The integers $x$ and $y$ are said to be
  \textbf{relatively prime} if and only if
  $\gcd(x,y) = 1$.
\end{define}\newpage


%%%%%%%%%%%%%%%%%%%%%%%%%%%%%%%%%%%%%%%%%%%%%%%%%%%%%%%%%%%%%%%%%%%%%%%
% Problems
%%%%%%%%%%%%%%%%%%%%%%%%%%%%%%%%%%%%%%%%%%%%%%%%%%%%%%%%%%%%%%%%%%%%%%%

\begin{enumerate} 
  \item Let $a \in \setZ$. Prove that $3\LDIV a^2$ if and only if $3\LDIV a$.
  \label{prob:divisibility-by-3} % Problem~\ref{prob:divisibility-by-3}
  \textit{You may use the fact that every integer can be written as 
  exactly one of $3k$, $3k + 1$, or $3k+2$ for some integer $k$.}
  \begin{subprove}
    Let $a \in \setZ$ be arbitrary. \vspace*{10px}\\
    ($\Rightarrow$) The contrapositive of the given statement is \inlinequote{$3\NDIV a$ only if $3\NDIV a^2$.}
    Assume $3\NDIV a$ for the sake of contrapositive proof. 
    By the given fact, $a$ can be written as exactly one of $3k + 1$ or $3k+2$ for some 
    integer $k$ since $a$ cannot be multiple of $3$.
    \begin{itemize}
      \item If $a = 3k + 1$, then $a^2 = (3k + 1)^2 = 9k^2 + 6k + 1 = 3(3k^2 + 2k) + 1$. 
      Since $3k^2 + 2k$ is also an integer by Closure Property of Integers, by definition of 
      divisibility, $3\NDIV a^2$.
      \item If $a = 3k + 2$, then $a^2 = (3k + 2)^2 = 9k^2 + 12k + 4 = 3(3k^2 + 4k + 1) + 1$. 
      Since $3k^2 + 4k + 1$ is also an integer by Closure Property of Integers, by definition of 
      divisibility, $3\NDIV a^2$.
    \end{itemize}
    In both cases, $3\NDIV a^2$.
    Therefore, by contrapositive, $3\LDIV a^2$ only if $3\LDIV a$.\vspace*{10px}\\
    ($\Leftarrow$) Assume $3\LDIV a$ for the sake of direct proof. 
    Since $3\LDIV a$, there exists an integer $k$ such that $a = 3k$ by definition of divisibility.
    Then $a^2 = (3k)^2 = 9k^2 = 3(3k^2)$.
    Since $3k^2$ is also an integer by Closure Property of Integers, by definition of divisibility, $3\LDIV a^2$.
    Thus, $3\LDIV a^2$ if $3\LDIV a$.\vspace*{10px}\\
    Therefore, by ($\Rightarrow$) and ($\Leftarrow$), $3\LDIV a^2$ if and only if $3\LDIV a$.
  \end{subprove}
  \newpage
  
  \item It is to be proven that $\sqrt{3}$ is irrational.
  \label{prob:sqrt3-irrational} % Problem~\ref{prob:sqrt3-irrational}
  \textit{Hint: The result from Problem~\ref{prob:divisibility-by-3} should be used.}
  \begin{subprove}
    Suppose $\sqrt{3}$ is rational for the sake of contradiction. By the definition of 
    a rational number, there exist integers $p$ and $q$, with $q \neq 0$, such that 
    $$\sqrt{3} = \frac{p}{q}.$$ 
    It is further assumed that the fraction $\frac{p}{q}$ has 
    been reduced to its lowest terms, meaning that $p$ and $q$ are relatively prime, 
    i.e., $\gcd(p,q) = 1$. Then, the equation will be rearranged as follows.
    $$p^2 = 3q^2.$$
    Since $q$ is an integer, $q^2$ is an integer by Closure Property of 
    Integers, which implies $3\LDIV p^2$ and $3\LDIV p$ by Problem~\ref{prob:divisibility-by-3}.
    Furthermore, by the definition of divisibility, there exists an integer $k$ such that $p = 3k$. 
    Thus,
    $$
    p^2 =(3k)^2 = 9k^2 = 3q^2.
    $$
    Moreover,
    $$
    3k^2 = q^2
    $$
    which implies again $3\LDIV q^2$ and $3\LDIV q$ by Problem~\ref{prob:divisibility-by-3}.
    This means that $p$ and $q$ share a common factor of $3$, which contradicts the assumption 
    that $\gcd(p,q) = 1$ (or $p$ and $q$ are relatively prime.) 
    Therefore, $\sqrt{3}$ is irrational.
  \end{subprove}
  \newpage


  \item Prove the following directly from the definition.   \\
  Let $a$, $b$, $c$, $d$, $x$ and $y$ be integers with $a \neq 0$ and $b \neq 0$.
  \begin{enumerate}
    \item If $a\LDIV c$ and $b\LDIV d$, then $ab\LDIV cd$.
    \begin{subprove}
      Assume $a\LDIV c$ and $b\LDIV d$ for the sake of direct proof.
      By definition of divisibility, there exist integers $k$ and $l$ such that $c = ak$ and $d = bl$.
      Then, $cd = (ak)(bl) = (ab)(kl)$. Since $kl$ is also an integer, by Closure Property 
      of Integers and definition of divisibility, $ab\LDIV cd$.
    \end{subprove}
    
    \item If $a\LDIV c$ and $a\LDIV d$, then $a\LDIV cx + dy$.
    \begin{subprove}
      Assume $a\LDIV c$ and $a\LDIV d$ for the sake of direct proof.
      By definition of divisibility, there exist integers $k$ and $l$ such that $c = ak$ and $d = al$.
      Then $cx + dy = (ak)x + (al)y = a(kx + ly)$.
      Since $kx + ly$ is an integer, by Closure Property 
      of Integers and definition of divisibility, $a\LDIV cx + dy$.
    \end{subprove}
    
    \item If $a\LDIV b$ and $b\LDIV a$, then $a = b$ or $a = -b$.
    \begin{subprove}
      Assume $a\LDIV b$ and $b\LDIV a$ for the sake of direct proof.
      By definition of divisibility, there exist integers $k$ and $l$ such that $b = ak$ and $a = bl$. 
      Without loss of generality of arguing $a$ and $b$ seperately, $b = (bl)k = b(lk)$ by substitution.
      Since $b \neq 0$ (and $a \neq 0$ as well) by condition, $1 = lk$ by dividing $b$ on both sides. The only integer solutions 
      to this equation are $(l,k) = (1,1)$ or $(l,k) = (-1,-1)$.
      If $(l,k) = (1,1)$, then $a = b$.
      If $(l,k) = (-1,-1)$, then $a = -b$.
      Therefore, $a = b$ or $a = -b$.
    \end{subprove}
    
    \item If $a\NDIV cd$, then $a\NDIV c$ and $a\NDIV d$.
    \begin{subprove}
      The contrapositive of the given statement is \inlinequote{if $a\LDIV c$ or $a\LDIV d$, then $a\LDIV cd$.}
      Assume $a\LDIV c$ or $a\LDIV d$ for the sake of direct contrapositive.
      \begin{itemize}
        \item If $a\LDIV c$, then there exists an integer $k$ such that $c = ak$. 
        Then $cd = (ak)d = a(kd)$. Since $kd$ is also an integer, $a\LDIV cd$.
        \item If $a\LDIV d$, then there exists an integer $l$ such that $d = al$. 
        Then $cd = c(al) = a(cl)$. Since $cl$ is also an integer, $a\LDIV cd$.
      \end{itemize}
      In both cases, $a\LDIV cd$. Therefore, by contrapositive, 
      if $a\NDIV cd$, then $a\NDIV c$ and $a\NDIV d$.
    \end{subprove}
  \end{enumerate}

  \newpage


  \item Let $A$, $B$ be subsets of a universal set $\setUNIV$.
  Prove the following using the element method:
  \begin{enumerate}
    \item $A \sMINUS B \subseteq A$.
    \begin{subprove}
      Let $x \in A \sMINUS B$ be arbitrary, for the sake of direct proof.
      By definition of set difference, $x \in A$ and $x \notin B$. Thus, $x \in A$ (trivial).
      Therefore, $A \sMINUS B \subseteq A$ by definition of Subset.
    \end{subprove}
    
    \item $A \sMINUS B = A \sAND \sNOT{B}$
    \begin{subprove}
      By definition, the equality $A \sMINUS B = A \sAND \sNOT{B}$ 
      requires showing both directions of subset relations, 
      $A \sMINUS B \subseteq A \sAND \sNOT{B}$ and 
      $A \sMINUS B \supseteq A \sAND \sNOT{B}$.\vspace*{10px}\\
      ($\subseteq$) Let $x \in A \sMINUS B$ be arbitrary.
      By definition of set difference, $x \in A$ and $x \notin B$.
      Since $x \notin B$ by definition of complementary set, it follows that $x \in \sNOT{B}$.
      Therefore, $x \in A$ and $x \in \sNOT{B}$, so $x \in A \sAND \sNOT{B}$.
      Hence, $A \sMINUS B \subseteq A \sAND \sNOT{B}$.\vspace*{10px}\\
      ($\supseteq$) Let $x \in A \sAND \sNOT{B}$ be arbitrary.
      By definition of intersection, $x \in A$ and $x \in \sNOT{B}$.
      Since $x \in \sNOT{B}$, it follows that $x \notin B$ by definition of complementary set.
      Therefore, $x \in A$ and $x \notin B$, so $x \in A \sMINUS B$.
      Hence, $A \sAND \sNOT{B} \subseteq A \sMINUS B$.\vspace*{10px}\\
      Therefore, by ($\subseteq$) and ($\supseteq$), $A \sMINUS B = A \sAND \sNOT{B}$.
    \end{subprove}
    
    \item $A \sAND (B \sMINUS A) = \setNULL$.
    \begin{subprove}
      Assume there exists an element $x \in A \sAND (B \sMINUS A)$ by negation of empty set definition for 
      the sake of contradiction. Then $x \in A$ and $x \in B \sMINUS A$.
      Since $x \in B \sMINUS A$, it follows that $x \in B$ and $x \notin A$, 
      which contradicts the assumption, $x \in A$.
      Therefore, $A \sAND (B \sMINUS A) = \setNULL$.
    \end{subprove}
    
    \item $A \sOR (B \sMINUS A) = A \sOR B$.
    \begin{subprove}
      By definition, the equality $A \sOR (B \sMINUS A) = A \sOR B$ 
      requires showing both directions of subset relations, 
      $A \sOR (B \sMINUS A) \subseteq A \sOR B$ and 
      $A \sOR (B \sMINUS A) \supseteq A \sOR B$.\vspace*{10px}\\
      ($\subseteq$) Let $x \in A \sOR (B \sMINUS A)$ be arbitrary.
      Then $x \in A$ or $x \in B \sMINUS A$.
      If $x \in A$, then $x \in A \sOR B$ as well (trivially).
      If $x \in B \sMINUS A$, then $x \in B$ and $x \notin A$, so $x \in B$, 
      which also (trivially) implies $x \in A \sOR B$.
      In both cases, $x \in A \sOR B$. Hence, $A \sOR (B \sMINUS A) \subseteq A \sOR B$.\vspace*{10px}\\
      ($\supseteq$) Let $x \in A \sOR B$ be arbitrary.
      Then $x \in A$ or $x \in B$. If $x \in A$, then $x \in A \sOR (B \sMINUS A)$ as well (trivially). 
      Next, if $x \in B$, then either $x \in A$ or $x \notin A$.
      \begin{itemize}
        \item If $x \in A$, then $x \in A \sOR (B \sMINUS A)$. (trivial)
        \item If $x \notin A$, then $x \in B \sMINUS A$ by definition of set 
        difference, so $x \in A \sOR (B \sMINUS A)$. (trivial)
      \end{itemize}
      In all cases, $x \in A \sOR (B \sMINUS A)$. Hence, $A \sOR B \subseteq A \sOR (B \sMINUS A)$.\vspace*{10px}\\
      Therefore, by ($\subseteq$) and ($\supseteq$), $A \sOR (B \sMINUS A) = A \sOR B$.
    \end{subprove}
  \end{enumerate}

  \newpage


  \item Let $A$, $B$, and $C$ be subsets of a universal set $\setUNIV$.
  Prove or disprove: 
  \begin{enumerate}
    \item If \(A \sAND B = A \sAND C\), then \(B = C\).
    \begin{subdisprove}
      Let $\setUNIV = \{1, 2, 3, 4\}$, $A = \{1, 2\}$, 
      $B = \{1, 3\}$, and $C = \{1, 4\}$ be subsets of a universal set $\setUNIV$
      for the sake of counterexample.
      Then, $A \sAND B = \{1\}$ and $A \sAND C = \{1\}$, that is, $A \sAND B = A \sAND C$.
      However, $B = \{1, 3\} \neq \{1, 4\} = C$. Therefore, it is \textbf{false} 
      that \(A \sAND B = A \sAND C\) only if \(B = C\).
    \end{subdisprove}
    
    \item If \(A \sOR B = A \sOR C\), then \(B = C\).
    \begin{subdisprove}
      Let $\setUNIV = \{1, 2, 3, 4\}$, $A = \{1, 2\}$, $B = \{2, 3\}$, and $C = \{3\}$ 
      be subsets of a universal set $\setUNIV$
      for the sake of counterexample.
      Then $A \sOR B = \{1, 2, 3\}$ and $A \sOR C = \{1, 2, 3\}$, so $A \sOR B = A \sOR C$.
      However, $B = \{2, 3\} \neq \{3\} = C$.
      Therefore, it is \textbf{false} that \(A \sOR B = A \sOR C\) only if \(B = C\).
    \end{subdisprove}
    
    \item If \(A \sAND B = A \sAND C\) and \(A \sOR B = A \sOR C\), then \(B = C\).
    \begin{subprove}
      Assume $A \sAND B = A \sAND C$ and $A \sOR B = A \sOR C$ for 
      the sake of direct proof.
      By definition of set equality, both directions of subset relations are 
      required as follows.\vspace*{10px}\\
      ($\subseteq$) Let $x \in B$ be arbitrary.
      Since $x \in B$, it follows that $x \in A \sOR B = A \sOR C$.
      Therefore, $x \in A$ or $x \in C$ (trivial). If $x \in A$, then 
      $x \in A \sAND B = A \sAND C$, so $x \in C$.
      Hence, $B \subseteq C$.\vspace*{10px}\\
      ($\supseteq$) Let $x \in C$ be arbitrary.
      Since $x \in C$, it follows that $x \in A \sOR C = A \sOR B$.
      Therefore, $x \in A$ or $x \in B$ (trivial).
      If $x \in A$, then $x \in A \sAND C = A \sAND B$, so $x \in B$.
      Hence, $C \subseteq B$.\vspace*{10px}\\
      Therefore, by ($\subseteq$) and ($\supseteq$), \(B = C\) when \(A \sAND B = A \sAND C\) and \(A \sOR B = A \sOR C\).
    \end{subprove}
  \end{enumerate}

  \newpage


  \item Let $A$, $B$ be subsets of a universal set $\setUNIV$.
  Prove or disprove:
  \begin{enumerate}
    \item $A \sTIMES B = \setNULL$ if $A = \setNULL$ or $B = \setNULL$.
    \begin{subprove}
      Assume $A = \setNULL$ or $B = \setNULL$ for the sake of direct proof.
      If $A = \setNULL$, then $A \sTIMES B = \setNULL \sTIMES B = \setNULL$ (since there are no first coordinate value, meaning there are no coordinates).
      If $B = \setNULL$, then $A \sTIMES B = A \sTIMES \setNULL = \setNULL$ (since there are no second coordinate value, meaning there are no coordinates).
      Therefore, direct proof by cases, if $A = \setNULL$ or $B = \setNULL$, then $A \sTIMES B = \setNULL$.
    \end{subprove}
    
    \item $A \sTIMES B = \setNULL$ only if $A = \setNULL$ or $B = \setNULL$.
    \begin{subprove}
      The contrapositive of the given statement is \inlinequote{if $A \neq \setNULL$ 
      and $B \neq \setNULL$, then $A \sTIMES B \neq \setNULL$.}
      Assume $A \neq \setNULL$ and $B \neq \setNULL$.
      Since $A \neq \setNULL$, there exists an element $a \in A$.
      Without loss of generality, there exists an element $b \in B$.
      Then $(a,b) \in A \sTIMES B$ by definition of Cartesian product.
      Therefore, $A \sTIMES B \neq \setNULL$ since there exists an element.
      Therefore, by contrapositive, if $A \sTIMES B = \setNULL$, then $A = \setNULL$ or $B = \setNULL$.
    \end{subprove}
  \end{enumerate}

\end{enumerate}